\chapter{Advanced Memory: Coherence and Consistency}\label{ch:memory-advanced}
\begin{center}\small\textit{Many caches, one memory: keep them coherent without strangling performance, and define what ``correct'' sharing means.}\end{center}

\section{Why Coherence Matters}
A chip has multiple cores, each with private L1/L2 caches sharing an address space. Without intervention, two cores can hold different values for the same address, violating the programmer's expectation of a single memory. \textbf{Coherence} ensures all cores eventually agree on each location's value; \textbf{consistency} defines when different locations' updates become visible and in what order.

Example failure without coherence:
\begin{center}\begin{tabular}{lcc}
\toprule Step & Core 0 & Core 1 \\ \midrule
0 & \texttt{x}=0 in DRAM, both caches empty & \\
1 & write \texttt{x=1}, hits L1 & reads \texttt{x} $\to$ 0 from DRAM (stale) \\
2 & evicts dirty line later & still holds stale 0 \\ \bottomrule
\end{tabular}\end{center}
Coherence protocols track who has a copy and who may write; consistency models restrict how cores may reorder loads and stores.

\section{Coherence Properties}
For a single address $A$, coherence enforces: (1) write propagation (a write eventually visible to all), (2) write serialization (all cores see writes to $A$ in the same order), and (3) read returns the latest serialized write (before any later write overwrites it). Mechanisms are \textbf{snooping} (bus broadcast, all caches listen) or \textbf{directory} (a home node tracks sharers and forwards requests).

Write policy interacts: write-back caches keep dirty copies and must write back on eviction or when another core wants the data; write-through would simplify but costs bandwidth.

\begin{center}
\begin{tabular}{llp{7cm}}
\toprule Message & Source & Effect \\ \midrule
\texttt{BusRd} & cache on read miss & directory / snoop asks sharers; supplier provides data \\
\texttt{BusRdX} & cache on write miss & read-exclusive + invalidate sharers \\
\texttt{Invalidate} & writer in S$\to$M & sharers move S$\to$I, send Ack \\
\texttt{Ack / Data} & sharer / memory & completes the transition, may supply dirty data \\ \bottomrule
\end{tabular}
\end{center}

A cache miss stalls the core (or at least that load) until the coherence transaction completes. Modern cores allow hits under miss (non-blocking caches, MSHRs) so other accesses to different lines proceed, but dependences on the missing line still stall.

\begin{definition}[Stable vs.\ transient states]
Stable states are M/E/S/I as seen between transactions. Internally a line passes through transient states (e.g., \texttt{IS} waiting for data, \texttt{SM} waiting for acks) while messages are in flight. Verification must check that transient states never deadlock and that write serialization is preserved despite overlapping transactions.
\end{definition}

\section{MSI Protocol}
The base three-state protocol attaches to each cache line one of: \textbf{M}odified (exclusive dirty, must write back), \textbf{S}hared (may be replicated, clean), \textbf{I}nvalid (no valid copy).

\begin{center}
\begin{tikzpicture}[font=\scriptsize, scale=0.92, auto]
\node[draw, fill=red!14, rounded corners, minimum width=1.8cm, minimum height=0.9cm] (I) at (0,0) {\textbf{I}};
\node[draw, fill=green!14, rounded corners, minimum width=1.8cm, minimum height=0.9cm] (S) at (5,0) {\textbf{S}};
\node[draw, fill=orange!16, rounded corners, minimum width=1.8cm, minimum height=0.9cm] (M) at (5,2.2) {\textbf{M}};
\draw[-{Stealth}, thick, blue!60!black] (I) -- (S) node[midway, below, font=\tiny] {Rd miss, others have it};
\draw[-{Stealth}, thick, blue!60!black, bend left] (I) to node[midway, right, font=\tiny] {Rd miss, no sharer} (M);
\draw[-{Stealth}, thick, green!50!black] (S) -- (I) node[midway, right, font=\tiny] {Invalidate / evict};
\draw[-{Stealth}, thick, violet!70!black] (S) -- (M) node[midway, right, font=\tiny] {Wr hit + invalidate};
\draw[-{Stealth}, thick, red!60!black] (M) -- (I) node[midway, left, font=\tiny] {Wr miss / evict + writeback};
\draw[-{Stealth}, thick, orange!70!black] (M) to[bend left=20] node[midway, above, font=\tiny] {Rd from other: supply + $\to$S} (S);
\draw[-{Stealth}, thick, gray, loop left] (S) to node[left, font=\tiny] {Rd hit} (S);
\draw[-{Stealth}, thick, gray, loop above] (M) to node[above, font=\tiny] {Rd/Wr hit} (M);
\node[below=0.7cm of I, font=\tiny, align=center, text=gray!60!black] {Invalid: no copy\\must fetch};
\end{tikzpicture}
\end{center}

Events from the core (Rd, Wr) and from the bus (BusRd, BusRdX, Invalidate) drive transitions. On a write miss (\texttt{Wr miss}) the cache issues \texttt{BusRdX} (read-exclusive), invalidates sharers, and enters M. On a read miss where another core holds M, that core supplies the data and both end in S (or the supplier writes back).

MSI allows sharing for reads but a write requires invalidating all sharers first, costing bus traffic on every write to shared data.

\subsection{MSI Walk-Through}
Initial: DRAM x=0, both caches I.

\begin{center}
\begin{tabular}{cllccc}
\toprule Step & Core action & Bus action & Core0 & Core1 & Data from \\ \midrule
1 & C0 Rd x & BusRd & S & I & DRAM \\
2 & C1 Rd x & BusRd & S & S & LLC / C0 \\
3 & C0 Wr x=1 & BusRdX (invalidate) & M & I & --- \\
4 & C1 Rd x & BusRd & S (downgrade, supply) & S & C0 cache-to-cache \\ \bottomrule
\end{tabular}
\end{center}
At step 3 Core0 must collect ack from Core1 before completing; the store stalls in the store buffer until acks arrive. At step 4 Core0 supplies the dirty line and writes back; both end S with the new value. Repeating this with MESI would keep the line in E after step 1 (no sharer), so step 3 would be silent (E$\to$M, no BusRdX).

\section{MESI: Adding Exclusive}
MSI cannot distinguish ``shared clean with one copy'' from ``shared clean with many copies'': a core that brings a line from DRAM to an empty system still enters S and must issue an invalidate on its first write, even though no sharer exists. MESI adds \textbf{E}xclusive (clean, sole copy): a line fetched with no other sharer enters E and can be silently upgraded to M on a write without bus traffic.

\begin{center}
\begin{tikzpicture}[font=\scriptsize, scale=0.90, auto]
\node[draw, fill=red!12, rounded corners, minimum width=1.6cm, minimum height=0.85cm] (I2) at (0,0) {I};
\node[draw, fill=yellow!15, rounded corners, minimum width=1.6cm, minimum height=0.85cm] (E) at (3.0,0) {E};
\node[draw, fill=green!14, rounded corners, minimum width=1.6cm, minimum height=0.85cm] (S2) at (6.0,0) {S};
\node[draw, fill=orange!18, rounded corners, minimum width=1.6cm, minimum height=0.85cm] (M2) at (3.0,2.2) {M};
\draw[-{Stealth}, thick, blue!60!black] (I2) -- (E) node[midway, below, font=\tiny] {Rd miss, no sharer};
\draw[-{Stealth}, thick, blue!60!black] (I2) to[bend left=18] node[midway, above, font=\tiny] {Rd miss, sharer exists} (S2);
\draw[-{Stealth}, thick, violet!70!black] (E) -- (S2) node[midway, below, font=\tiny] {BusRd from other};
\draw[-{Stealth}, thick, green!60!black] (E) -- (M2) node[midway, left, font=\tiny] {Wr hit (silent)};
\draw[-{Stealth}, thick, orange!70!black] (S2) -- (M2) node[midway, right, font=\tiny] {Wr hit + inval};
\draw[-{Stealth}, thick, red!60!black] (S2) -- (I2) node[midway, below, font=\tiny] {Invalidate/evict};
\draw[-{Stealth}, thick, red!60!black] (M2) -- (I2) node[midway, left, font=\tiny] {evict + wb / BusRdX};
\draw[-{Stealth}, thick, gray, loop below] (S2) to node[below, font=\tiny] {Rd hit} (S2);
\draw[-{Stealth}, thick, gray, loop above] (M2) to node[above, font=\tiny] {Rd/Wr hit} (M2);
\end{tikzpicture}
\end{center}

E is the performance win: private data and stack lines stay in E/M without coherence traffic. Transitions: E$\to$S on snooped BusRd (now shared), E$\to$M on \texttt{Wr hit} silently (no invalidate needed), M$\to$I on eviction or when another core's \texttt{BusRdX} requests it (supply data and write back if needed).

MOESI (Adds \textbf{O}wned: dirty but shared, supplier without writeback) and MESIF (Adds \textbf{F}orward: designated responder) further optimise sharing of dirty lines and reduce broadcasts on read misses to shared data. Directory protocols scale similarly but the directory replaces bus snooping with point-to-point messages.

\section{Snooping versus Directory}
Snooping broadcasts every coherence request on a bus or ring; every cache checks its tags. Simple but bandwidth-limited to $\approx 8$--$16$ cores. Directory keeps a per-line sharer vector at the home node (often LLC/ memory controller); requests go to the directory, which forwards only to sharers. Directory scales to hundreds of cores but adds an indirection hop and directory storage overhead.

\begin{center}
\begin{tikzpicture}[font=\scriptsize, scale=0.88]
\node[draw, fill=blue!14, minimum width=1.8cm, minimum height=0.7cm] (c0) at (0,1.5) {Core 0 / L1};
\node[draw, fill=green!14, minimum width=1.8cm, minimum height=0.7cm] (c1) at (4,1.5) {Core 1 / L1};
\node[draw, fill=orange!16, minimum width=2.6cm, minimum height=0.8cm] (dir) at (2,0) {Directory / LLC};
\node[draw, fill=gray!16, minimum width=1.8cm, minimum height=0.7cm] (mem) at (2,-1.2) {DRAM};
\draw[{Stealth}-{Stealth}, thick, blue!60!black] (c0.south) -- (dir.north west) node[midway, left, font=\tiny] {request/forward};
\draw[{Stealth}-{Stealth}, thick, green!60!black] (c1.south) -- (dir.north east);
\draw[{Stealth}-{Stealth}, thick, gray] (dir.south) -- (mem.north);
\node[right=0.4cm of c1, font=\tiny, text=gray!60!black] {snoop bus also shown as arrows via directory};
\end{tikzpicture}
\end{center}

\section{False Sharing, Granularity and MOESI Extensions}
Coherence acts on whole lines (64\,B). Two independent variables on the same line cause \textbf{false sharing}: cores ping-pong the line between M states even though they touch disjoint bytes. Example: \texttt{int cnt[2]} with \texttt{cnt[0]} on Core0 and \texttt{cnt[1]} on Core1 still bounces every write because the line is one coherence unit. Fixes: pad structures to line boundaries, align thread-local data, or use per-core partitioning.

\begin{center}
\begin{tabular}{lcccc}
\toprule Extension & Extra state & Meaning & When it helps \\ \midrule
MOESI & O (Owned) & dirty but shared, owner supplies without writeback & producer--consumer sharing \\
MESIF & F (Forward) & designated sharer to reply & many readers, one supplier \\
MESI & E (Exclusive) & clean sole copy & private/stack data \\ \bottomrule
\end{tabular}
\end{center}

Owned (MOESI, AMD) lets a dirty line be shared without writing back: the owner (O) supplies data on snooped reads and remembers to write back on eviction. Forward (MESIF, Intel) picks one sharer as responder so multiple S copies do not all try to respond. Neither changes correctness, only traffic. Directory with MOESI can keep O in the LLC and serve reads without going to the owning core's L1.

Coherence granularity also affects bandwidth: a write to one byte still invalidates the whole line. Sector caches and write-combining buffers mitigate by merging multiple writes to the same line before issuing a single invalidate.

\section{Coherence Performance}
Misses to shared lines cost more than private misses: a read miss to an S line may be served from LLC ($\approx 10$\,ns) but a read miss needing data from another core's M line costs a cache-to-cache transfer ($\approx 20$--$40$\,ns plus coherence messages). Invalidations are on the critical path for writes to shared data: the store cannot commit until all sharers acknowledge invalidation, which the LSQ and write buffer must track. Large sharing (many sharers) therefore stalls the writer longer.

A rule of thumb: minimise sharing, maximise private/E-state data, and align synchronisation variables (locks, flags) on their own lines away from payload data. Tools like false-sharing detectors sample coherence misses via performance counters (e.g., \texttt{MESI\_Snoop\_Hit}, \texttt{Offcore\_Response}) to guide padding.

\section{Memory Consistency}
Coherence says each location is serialized; consistency says how different locations' accesses interleave as seen by different cores. Consider initial \texttt{x=y=0}:

\begin{lstlisting}[language={[x86masm]Assembler}, caption={Classic litmus test (store buffering)}, label={lst:litmus1}]
# Core 0              # Core 1
SW  x1, 0(a0) # x=1    SW  x1, 0(a1) # y=1
LW  x2, 0(a1) # y=?    LW  x2, 0(a0) # x=?
# Under SC, at least one load sees 1. Under x86-TSO or weaker, both may see 0.
\end{lstlisting}

\textbf{Sequential Consistency (SC)}: there exists a global interleaving of all cores' operations that respects each core's program order and where each load sees the latest store in that interleaving. Intuitive but expensive: it forbids most OOO and buffering optimisations.

\textbf{Total Store Order (TSO, x86)}: loads may be reordered before earlier stores to different addresses (store buffer), but stores stay ordered with stores, loads with loads. The store-buffer litmus above is allowed to see (0,0); earlier litmus where both stores are seen is still forbidden.

\textbf{Weak / release consistency (ARM, RISC-V) }: almost any reordering is allowed except when the programmer inserts fences. RISC-V RVWMO defines \texttt{FENCE}, \texttt{acquire} and \texttt{release} annotations: a release store makes prior accesses visible before it, an acquire load makes later accesses start after it. Data-race-free programs with proper synchronisation appear SC; racy code has no guarantee.

\begin{lstlisting}[language={[x86masm]Assembler}, caption={Fences restore ordering in RISC-V}, label={lst:fence}]
# Core 0: producer                    # Core 1: consumer
SW  x1, 0(a0)      # data=1              LW  t0, 0(a1)  # flag ?
FENCE  w, w                            FENCE  r, r
SW  x1, 0(a1)      # flag=1              LW  t1, 0(a0)  # data
# Without fences on weak model, consumer could see flag=1 but data=0.
# FENCE w,w orders stores; FENCE r,r orders loads.
\end{lstlisting}

Fences flush or order store buffers and drain load queues; they cost tens of cycles, so place them only at synchronisation points. Compilers map C++ \texttt{std::atomic} with \texttt{memory\_order\_release/acquire} to exactly these fences.

\subsection{What Each Model Allows to Reorder}
\begin{center}
\begin{tabular}{lcccc}
\toprule Reordering & SC & TSO & ARMv8 & RVWMO \\ \midrule
Store $\to$ later Load (diff addr) & No & Yes (SB) & Yes & Yes \\
Store $\to$ later Store & No & No & Yes* & Yes* \\
Load $\to$ later Load & No & No & Yes* & Yes* \\
Load $\to$ later Store & No & No & Yes* & Yes* \\ \bottomrule
\end{tabular}
\end{center}
{\footnotesize *Allowed unless ordered by \texttt{FENCE} or \texttt{acquire/release} annotations. SC allows none; TSO only relaxes Store$\to$Load via store buffer; weak models relax all four.}

The classic store-buffer allows core 0 to buffer \texttt{x=1} while loading \texttt{y} before the store is globally visible; both cores can then see (0,0). Under SC the store buffer must drain before the load, so at least one load sees 1. On RISC-V, a \texttt{FENCE RW,RW} between the store and load would restore SC for that pair. Acquire/release is lighter than a full fence: \texttt{AMOSWAP.aq} (acquire) orders later ops after it, \texttt{AMOSWAP.rl} (release) orders earlier ops before it, matching C++ atomics without a heavyweight \texttt{FENCE}.

A second litmus is \textit{message passing}: \texttt{data=1; flag=1} on Core0, \texttt{if(flag) load data} on Core1. Without \texttt{FENCE W,W} after data and \texttt{FENCE R,R} before the data load, flag=1 does not imply data=1 on a weak model. With fences, or with \texttt{SW.rl flag} and \texttt{LW.aq flag}, the ordering is enforced.

\section{Putting It Together: Coherence vs.\ Consistency}
Coherence is per-location and enforced by MSI/MESI; consistency is cross-location and enforced by fences and store-buffer draining. A core with aggressive OOO may reorder a load past an earlier store (TSO-allowed) but must not reorder two stores to the same address (coherence violation). Modern cores speculate past fences and roll back on violation detection via the LSQ and coherence invalidations.

Speculative coherence adds nuance: a core may speculatively read a line in S, but if another core invalidates it before the speculative load commits, the LSQ must detect the invalidation and replay. Store buffers also interact: a store sits in the buffer until its line is in M/E; only then does it become globally visible. Draining the buffer is the consistency action; invalidating sharers is the coherence action.

% ============================================================
\section{Cache Optimisations Consolidated: Six Classic Tricks}\label{sec:cache-opts}
Earlier chapters borrow one cache trick each in passing: out-of-order execution hiding a miss (Chapter~2), the fetch chain of I-cache, BTB and prefetch (Chapter~3), and coherence traffic for shared lines (Chapter~7). This section collects the six standard uniprocessor optimisations in one place.
\begin{center}
\begin{tabular}{lp{8.2cm}}
\toprule Trick & What it buys \\ \midrule
Victim cache & small fully associative buffer catches conflict-evicted lines \\
Stride/stream prefetcher & detects sequential or strided access, fetches ahead \\
Non-blocking (lockup-free) & MSHRs track misses so hits proceed underneath \\
Critical-word-first + early restart & requested word first, resume before line fills \\
TLB organisation & caches translations so hits avoid page walks \\
Sv39 interaction & 3-level RISC-V tables set the walk cost on a TLB miss \\ \bottomrule
\end{tabular}
\end{center}
\textbf{Victim cache} (Jouppi): a 4--16-entry fully associative buffer between L1 and the refill path holds just-evicted lines, turning conflict misses into fast victim hits without raising L1 associativity. \textbf{Hardware prefetchers} recognise streams and strides, such as the unit-stride sweeps of a stencil, and fetch lines before the core asks; accuracy matters because wrong-path prefetches waste bandwidth and pollute the cache. \textbf{Non-blocking (lockup-free) caches} attach Miss Status Holding Registers (MSHRs) to each outstanding miss so later hits to other lines proceed underneath, which is exactly what lets the Chapter~2 out-of-order window keep draining during a miss. \textbf{Critical-word-first} delivers the requested word over the bus before the rest of the line, and \textbf{early restart} resumes the stalled load as soon as that word arrives rather than waiting for the full 64\,B fill. The \textbf{TLB} caches virtual-to-physical translations so hits never walk the page table; on a miss under RISC-V \textbf{Sv39} (39-bit virtual address, three-level radix table) the walker may issue up to three dependent memory reads, each itself cacheable. Hence TLB reach (entries $\times$ page size, grown by superpages) must cover the working set, or page-walk traffic lands on the same interconnect analysed in Chapter~5.

\section{Key Takeaways}
MSI is the minimal coherent protocol; MESI's Exclusive state eliminates unnecessary invalidates for private data. Snooping versus directory is a scalability trade-off. Coherence is necessary but not sufficient: the consistency model (SC/TSO/RVWMO) specifies which cross-address reorderings the hardware may perform. Programmers rely on synchronisation (locks, fences, acquire/release) to restore intuitive ordering where it matters, paying fence cost only on those paths.

\begin{remark}[Checklist]
For any litmus test, ask: (1) is coherence violated (per-location order)? (2) is consistency violated (cross-location order)? (3) would a fence or an \texttt{.aq/.rl} fix it, and at what cost? Answering all three is the exam pattern for \texttt{flag/data} and store-buffer tests.
\end{remark}

\textit{Next}: with coherence and consistency in hand, Chapter~5 extends the hierarchy to LLC, interconnect and real SoC implications.

% SC3050 ch04 coherence and consistency complete.

\section{Exercises}
\begin{exercise}Two cores, line \texttt{A} initially I in both, DRAM=0. Trace MSI states cycle by cycle: Core0 Rd A, Core1 Rd A, Core0 Wr A, Core1 Rd A. List bus transactions (BusRd/BusRdX/invalidate), data supplier, and final states. Repeat for MESI and highlight where E avoids a transaction.\end{exercise}
\begin{exercise}For a snooping bus, explain why a write to a line in S must broadcast an invalidate even under TSO. What would break if the invalidation were deferred to the store buffer drain? Give a litmus test where deferral violates coherence.\end{exercise}
\begin{exercise}Store-buffer litmus from Listing~1: both cores see 0 under TSO but not under SC. Show a TSO execution with store buffers where both loads bypass their own buffered stores. Draw the global order that SC would require and why it cannot include (0,0).\end{exercise}
\begin{exercise}RISC-V weak model: thread 0 does \texttt{SW data=1; FENCE w,w; SW flag=1}, thread 1 does \texttt{LW flag; FENCE r,r; LW data}. Prove that if thread 1 sees \texttt{flag=1} it must see \texttt{data=1}. What happens without the fences? Give an interleaving where \texttt{flag=1, data=0} is visible.\end{exercise}
\begin{exercise}Directory vs.\ snooping: a 64-core chip runs a read-sharing workload (all cores read line \texttt{A} repeatedly). Compare bus traffic per read miss for snooping (broadcast to 63) vs.\ directory (ask directory, get data). Now consider a private-data workload (each core writes its own line). Which scales better and why does MESI E state matter more for private data?\end{exercise}
